\documentclass[a4paper,11pt]{ltjsarticle}
\usepackage[haranoaji]{luatexja-preset}
\usepackage{amsmath,amssymb}
\usepackage[top=12mm,bottom=10mm,left=14mm,right=14mm]{geometry}
\usepackage{array}
\usepackage{tikz}
\usetikzlibrary{calc}
\pagestyle{empty}
\setlength{\parindent}{0pt}
\newlength{\qcol}\setlength{\qcol}{0.47\textwidth}
\newlength{\qgap}
\newlength{\anscolw}\setlength{\anscolw}{0.30\textwidth}
\newlength{\ansh}
\newcommand{\Q}[2]{\parbox[t]{\qcol}{\raggedright (#1)\hspace{0.6em}$\displaystyle #2$}}
\newcommand{\QR}[4]{\Q{#1}{#2}\hfill\Q{#3}{#4}\par\vspace{\qgap}}
\newcommand{\anscell}[1]{\rule[-\ansdrop]{0pt}{\ansh}(#1)}
\newlength{\ansdrop}

\begin{document}
{\large\bfseries 計算問題　8}\hfill 名前(\underline{\hspace{32mm}})\par\vspace{3mm}
■（1）〜（16）の計算をしなさい。（17）〜（20）は方程式を解きなさい。\par\vspace{3mm}
\setlength{\qgap}{5.4mm}
\QR{1}{(5x+9)-(4x-7)}{2}{(3a+5)+(-2a-8)}
\QR{3}{54x\div(-6)}{4}{(-7)\times5a}
\QR{5}{2+(-3)^2+(-5)}{6}{(-36)\div(-4)}
\QR{7}{-5+(-6)}{8}{(-7)\times(-8)}
\QR{9}{\dfrac{x-7}{4}\times8}{10}{(60x-12)\div6}
\QR{11}{(+6)+(-8)-(-7)}{12}{9\div4\times(-20)}
\QR{13}{(-7)-(-10)}{14}{\dfrac{1}{4}(8x-12)+\dfrac{1}{3}(3x+6)}
\QR{15}{6-(-4)\div\left(-\dfrac{2}{5}\right)}{16}{\dfrac{1}{5}(10a+5)-\dfrac{1}{2}(8a-4)}
\QR{17}{3x+13=6x-5}{18}{2x-9=-x}
\QR{19}{0.3x-0.8=4}{20}{\dfrac{5}{4}x+1=\dfrac{4}{3}x}
\vspace{1mm}
\Q{21}{75\times16}\par\vspace{3mm}
\setlength{\ansh}{9.5mm}\setlength{\ansdrop}{6mm}%
\begin{center}\begin{tabular}{|p{\anscolw}|p{\anscolw}|p{\anscolw}|}\hline
\anscell{1} & \anscell{2} & \anscell{3} \\\hline
\anscell{4} & \anscell{5} & \anscell{6} \\\hline
\anscell{7} & \anscell{8} & \anscell{9} \\\hline
\anscell{10} & \anscell{11} & \anscell{12} \\\hline
\anscell{13} & \anscell{14} & \anscell{15} \\\hline
\anscell{16} & \anscell{17} & \anscell{18} \\\hline
\anscell{19} & \anscell{20} & \anscell{21} \\\hline
\end{tabular}\end{center}

\end{document}
